\chapter{KNOWN ISSUES AND NEXT STEPS}
\label{ch:issues}

\section{OPEN ISSUES}

\begin{description}
\item[Board resets under full load.] The cape's \SI{5}{V} rail drops out with
  the camera, motors, lidar and Bluetooth all running (\S\ref{sec:powerbudget}).
  Waiting on a USB splitter and battery bank to power the camera separately.
\item[SLAM loses track while driving.] Analysis in \S\ref{sec:slamdiag}. The
  parked drift is fixed on branch \texttt{slam-estimator-updates}; the lidar
  mounting yaw is the leading remaining suspect and needs a drive test.
\item[Skid-steer turning.] Low turn commands cannot overcome scrubbing friction
  (\S\ref{sec:skidturn}). The teleop curve is a workaround; closed-loop
  wheel-speed control would fix it properly.
\item[Bluetooth input stalls.] Controller reports pause for 100--500~ms in
  bursts. ROS itself was measured clean (no \topic{/cmd\_vel} gap over
  \SI{58}{ms}). Parked as acceptable for now.
\item[Placeholder static transforms.] The lidar, IMU and GPS offsets in
  \file{robot.launch} have not been measured.
\item[Constant odometry covariance.] \topic{/odom} reports a fixed pose
  covariance (\S\ref{sec:odom}). Harmless for the current EKF, which fuses
  only $v_x$, but wrong for any consumer of the pose.
\item[Repository behind the robot.] The compass, ZUPT, GPS fusion and gmapping
  changes are on \texttt{slam-estimator-updates}, not \texttt{main}. The
  teleop turn curve is only on the board.
\item[Chase turning will likely stall.] Chase never asks for more than about
  0.2 duty to turn, below what a spin on the floor needs
  (\S\ref{sec:chaselaw}).
\item[Detector field of view unverified.] The $300\times300$ preview is
  probably a centre crop of the colour frame, about \SI{42}{\degree} wide
  rather than \SI{69}{\degree}, which would make the chase fallback bearing
  too large (\S\ref{sec:previewcrop}).
\item[Chase stops on one missed frame.] A single missed detection resets the
  hit counter (\S\ref{sec:chasemeas}).
\item[No camera transform.] No static transform for the OAK-D-Lite is
  published.
\item[Pause and Mode buttons.] The cape's buttons still do not work.
\end{description}

\section{SUGGESTED IMPROVEMENTS}

In rough order of value for the current goals:
\begin{enumerate}
\item Verify and set the lidar's yaw in \frm{base\_link}$\to$\frm{laser}
  (\S\ref{sec:slamdiag}).
\item Map with gentle turns, under about \SI{0.5}{rad/s}
  (\S\ref{sec:skew}).
\item Close the speed loop and add a gyro yaw-rate loop
  (\S\ref{sec:speedloop}--\ref{sec:yawloop}). This makes turning independent
  of the floor and gives the chase controller the turn rate it asks for.
\item Bring chase up in the order of \S\ref{sec:chaseplan} once the camera
  has its own power, then add the target filter of \S\ref{sec:targetkf}.
\item Log \topic{/imu/data\_raw} at rest for an hour and fit the
  error model \eqref{eq:sensormodel}--\eqref{eq:walkingbias} with Allan
  variance, as in the thesis (thesis \S7.5). That gives principled process
  noise for the EKF and a better ZUPT time constant.
\item Publish a growing pose covariance from \eqref{eq:odomcov} instead of the
  constant one.
\item Fuse $v_y=0$ in the EKF and set a Mahalanobis rejection threshold on
  GPS (\S\ref{sec:ekfcfg}).
\item Consider slam\_toolbox, which handles loop closure and lifelong mapping
  better than gmapping and can switch to localisation on a saved map.
\item Align \frm{map} and \frm{gps\_map} with \eqref{eq:align} once there is
  outdoor data.
\end{enumerate}
